\\abstract{Hallucinations in Vision-Language Models (VLMs) --- where models generate plausible but visually ungrounded or factually incorrect content --- remain a significant barrier to reliable multimodal AI deployment. Existing mitigation strategies either require expensive model retraining or rely on static knowledge bases that cannot adapt to unseen visual contexts. We present CMVKG-Guard, a training-free middleware framework that intercepts and corrects hallucinations token-by-token during autoregressive generation, without modifying the underlying VLM. The framework dynamically constructs a per-inference Hierarchical Multimodal Knowledge Graph (DH-MMKG) from visual scene entities and enriches it via live queries to Wikidata and ConceptNet. Each candidate token is scored by a three-layer Cross-Modal Verification Engine (CMVE) combining visual-textual alignment, knowledge graph grounding, and multi-hop reasoning consistency; tokens falling below an adaptive, step-aware threshold trigger a Constrained Look-Ahead Beam Search that selects a grounded replacement and teacher-forces it into the decoding sequence. Experiments on six standard benchmarks --- including POPE, CHAIR, and MMHalBench --- against eight baselines show that CMVKG-Guard reduces object hallucination rates by 84\% on CHAIR, achieves 91.5\% F1 on POPE (outperforming the strongest training-free baseline by 5.8 percentage points), and incurs a median end-to-end latency of 12ms per token on a single A100 GPU. These results suggest that dynamic, inference-time knowledge grounding is a promising direction for hallucination mitigation in open-domain VLM settings.}
